% !TeX root = ../navier-stokes-manuscript.tex
\subsection{The local construction behind the maximal history}
\label{app:LocalConstruction}

One local theorem must do two jobs without changing its clock.  It constructs
the initial history and later restarts it from the terminal datum.  The clock is
set by \(\nu\) and the datum’s \(H^3\) size; any higher regularity already carried
by that datum survives on the same interval.

For real \(s\), let
\[
  H^s_\sigma(\R^3)
  :=\{v\in H^s(\R^3;\R^3):\nabla\cdot v=0\}.
\]

\begin{proposition}[Local construction and continuation alternative]\label[proposition]{prop:LocalMaximalHistory}
Fix \(t_0\geq0\), \(\nu>0\), and \(a\in H^3_\sigma(\R^3)\).  A time
\(\delta=\delta(\nu,\norm{a}_{H^3})>0\) exists for which there is a unique
normalized-pressure pair \((w,q)\) satisfying
\[
  w\in C([t_0,t_0+\delta];H^3_\sigma)
  \cap L^2(t_0,t_0+\delta;H^4),
  \qquad
  w_t\in L^2(t_0,t_0+\delta;H^2).
\]
If the datum also lies in \(H^m_\sigma\) for an integer \(m>3\), this same
interval carries
\[
  w\in C([t_0,t_0+\delta];H^m_\sigma)
  \cap L^2(t_0,t_0+\delta;H^{m+1}),
  \qquad
  w_t\in L^2(t_0,t_0+\delta;H^{m-1}).
\]
These local pairs agree on overlaps and form one maximal history.  If its
maximal time is finite, its \(H^3\) norm is unbounded along times approaching
that endpoint.
\end{proposition}

\begin{proof}[Proof of Proposition~\ref{prop:LocalMaximalHistory}]
Keep \(t_0,\nu,a\) fixed as above and measure all real Sobolev orders with
\[
  \norm{f}_{H^s}^2
  :=\int_{\R^3}\langle\xi\rangle^{2s}
       |\widehat f(\xi)|^2\,d\xi,
  \qquad
  \langle\xi\rangle=(1+|\xi|^2)^{1/2}.
\]
Write \(\Leray\) for the Leray projection and \(R_j\) for the Riesz transform
with symbol \(i\xi_j/|\xi|\).  Each acts boundedly on every \(H^s\).

\divider{The three rows forced by the nonlinearity and the heat equation}

The construction must first absorb one derivative of \(w\otimes w\).  In three
dimensions,
\[
  H^3(\R^3)\hookrightarrow W^{1,\infty}(\R^3),
  \qquad
  \norm{fg}_{H^3}\leq C\norm{f}_{H^3}\norm{g}_{H^3},
\]
so an \(H^3\) velocity is Lipschitz and products stay in \(H^3\).  After the
divergence and Leray projection this gives
\begin{equation}
  \norm{\Leray\nabla\!\cdot(f\otimes g)}_{H^2}
  \leq C\norm{f}_{H^3}\norm{g}_{H^3}.
  \label{eq:LocalProductEstimate}
\end{equation}
Leibniz’ rule proves the estimate for integer derivatives, using
\(H^2\hookrightarrow L^\infty\) and \(H^1\hookrightarrow L^6\); density then
removes the smoothness assumption.

Once the force lies in \(H^2\), the heat equation places one parabolic derivative
above the base state and the time derivative one below it.  Thus
\(z\in L^2_tH^4_x\) and \(z_t\in L^2_tH^2_x\), with \(H^3\) exactly between
them.  For any such \(v\),
\[
  v\in L^2(J;H^4),
  \qquad
  v_t\in L^2(J;H^2),
\]
frequency truncation followed by the fundamental theorem of calculus and
Cauchy--Schwarz yields
\begin{equation}
  \left|
    \norm{v(t)}_{H^3}^2-\norm{v(r)}_{H^3}^2
  \right|
  \leq
  2\norm{v_t}_{L^2(r,t;H^2)}
   \norm{v}_{L^2(r,t;H^4)}.
  \label{eq:LocalAdjacentTrace}
\end{equation}
The right-hand side vanishes with the length of the time interval.  Removing the
frequency cutoff gives a unique representative in \(C(J;H^3)\), so
the three rows of the strong class are
\[
  C_tH^3_x\cap L^2_tH^4_x,
  \qquad
  \partial_t v\in L^2_tH^2_x.
\]
The continuous middle trace is what lets the equation remember its initial
value while the heat operator gains one square-integrable derivative.

\divider{A linear estimate aligned with those rows}

On \(J_\delta=[t_0,t_0+\delta]\), \(0<\delta\leq1\), consider
\[
  z_t-\nu\Delta z=F,
  \qquad
  z(t_0)=a,
  \qquad
  F\in L^2(J_\delta;H^2).
\]
At the \(H^3\) level the force and solution occupy the adjacent spaces
\(H^2\) and \(H^4\).  Their Fourier pairing is
\[
  \langle F,z\rangle_{H^2,H^4;H^3}
  :=\int_{\R^3}\langle\xi\rangle^6
       \widehat F(\xi)\cdot\overline{\widehat z(\xi)}\,d\xi.
\]
The weight factors as
\(
  \langle\xi\rangle^6
  =\langle\xi\rangle^2\langle\xi\rangle^4
\)
and Cauchy--Schwarz gives
\begin{equation}
  \left|\langle F,z\rangle_{H^2,H^4;H^3}\right|
  \leq\norm{F}_{H^2}\norm{z}_{H^4}.
  \label{eq:LocalShiftedPairing}
\end{equation}
For a smooth solution the corresponding energy row is
\[
  \frac12\frac d{dt}\norm{z}_{H^3}^2
  +\nu\norm{\nabla z}_{H^3}^2
  =\operatorname{Re}\langle F,z\rangle_{H^2,H^4;H^3}.
\]
The inhomogeneous norm separates the dissipative part exactly:
\[
  \norm{z}_{H^4}^2
  =\norm{z}_{H^3}^2+\norm{\nabla z}_{H^3}^2.
\]
Insert \eqref{eq:LocalShiftedPairing}, absorb by Young’s inequality, and apply
Gronwall to obtain
\[
  \norm{z}_{C(J_\delta;H^3)}
  +\norm{z}_{L^2(J_\delta;H^4)}
  \leq
  C_\nu\left(
    \norm{a}_{H^3}+\norm{F}_{L^2(J_\delta;H^2)}
  \right).
\]
The equation then controls the missing time row as well:
\begin{equation}
  \norm{z}_{C(J_\delta;H^3)}
  +\norm{z}_{L^2(J_\delta;H^4)}
  +\norm{z_t}_{L^2(J_\delta;H^2)}
  \leq
  C_\nu\left(
    \norm{a}_{H^3}+\norm{F}_{L^2(J_\delta;H^2)}
  \right).
  \label{eq:LocalLinearEstimate}
\end{equation}
No constant here deteriorates as \(\delta\) shrinks within \((0,1]\).
Frequency-truncate \(a\) and \(F\), solve the smooth equations, and apply
\eqref{eq:LocalLinearEstimate} to differences.  The Cauchy limit solves the
untruncated problem, while \eqref{eq:LocalAdjacentTrace} recovers its prescribed
strong \(H^3\) value at \(t_0\).  Applying the estimate to the difference of two
solutions proves uniqueness.

\divider{Closing the nonlinear map}

Set
\[
  X_\delta
  :=C(J_\delta;H^3_\sigma)\cap L^2(J_\delta;H^4),
  \qquad
  \norm{w}_{X_\delta}
  :=\norm{w}_{C H^3}+\norm{w}_{L^2H^4}.
\]
The projected equation defines the map
\begin{equation}
  \Phi(w)(t)
  :=e^{\nu(t-t_0)\Delta}a
  -\int_{t_0}^{t}e^{\nu(t-s)\Delta}
    \Leray\nabla\!\cdot(w\otimes w)(s)\,ds.
  \label{eq:LocalMildMap}
\end{equation}
All abbreviated time norms below refer to \(J_\delta\).  The product estimate
and the factor \(\delta^{1/2}\) from the interval give
\begin{align}
  \norm{\Leray\nabla\!\cdot(w\otimes w)}_{L^2H^2}
  &\leq C\delta^{1/2}\norm{w}_{C H^3}^2,
  \label{eq:LocalQuadraticCount}
  \\
  \norm{\Leray\nabla\!\cdot(w\otimes w-v\otimes v)}_{L^2H^2}
  &\leq C\delta^{1/2}
  \bigl(\norm{w}_{C H^3}+\norm{v}_{C H^3}\bigr)
  \norm{w-v}_{C H^3}.
  \label{eq:LocalDifferenceCount}
\end{align}
The linear estimate turns these force bounds into
\[
  \norm{\Phi(w)}_{X_\delta}
  \leq C_\nu\norm{a}_{H^3}
  +C_\nu C\delta^{1/2}\norm{w}_{X_\delta}^2
\]
and
\[
  \norm{\Phi(w)-\Phi(v)}_{X_\delta}
  \leq C_\nu C\delta^{1/2}
  \bigl(\norm{w}_{X_\delta}+\norm{v}_{X_\delta}\bigr)
  \norm{w-v}_{X_\delta}.
\]
When \(a\neq0\), set \(R=2C_\nu\norm{a}_{H^3}\) and choose
\(\delta\leq1\) with \(4C_\nu C\delta^{1/2}R\leq1\).  Then \(\Phi\) maps the
closed radius-\(R\) ball into itself and contracts it by at least a factor of
two.  For \(a=0\), the fixed point is identically zero.  Banach’s theorem gives
\(w\in X_\delta\), with a choice satisfying
\begin{equation}
  \delta
  \geq
  \min\left\{1,c_\nu(1+\norm{a}_{H^3})^{-2}\right\}.
  \label{eq:LocalLifespan}
\end{equation}
The constant \(c_\nu>0\) depends only on the fixed viscosity and the chosen
Sobolev normalization.

Uniqueness is not confined to that ball.  If \(w,v\in X_\delta\) share the
datum, apply the linear and difference estimates to \(w-v\) on an initial
subinterval short enough to absorb the coefficient.  The solutions agree
there.  Their common endpoint starts the same argument on the next subinterval,
and a finite partition covers \(J_\delta\).  The heat semigroup and \(\Leray\)
preserve divergence-free fields; the projected equation also supplies
\begin{equation}
  w_t\in L^2(J_\delta;H^2).
  \label{eq:LocalTimeRow}
\end{equation}

\divider{Restoring the pressure to the same solution}

For the constructed velocity, set
\begin{equation}
  q:=R_iR_j(w_iw_j).
  \label{eq:LocalPressureFormula}
\end{equation}
The algebra property and Riesz bounds imply
\[
  q\in C(J_\delta;H^3).
\]
For \(F=\nabla\!\cdot(w\otimes w)\), the Fourier symbols identify
\(\Leray F=F+\nabla q\).  Hence the projected fixed point is the
pressure-complete solution
\[
  w_t+(w\cdot\nabla)w+\nabla q=\nu\Delta w,
  \qquad
  \nabla\cdot w=0,
  \qquad
  w(t_0)=a.
\]
Membership of \(q(t)\) in \(H^3(\R^3)\) removes the whole-space additive
constant.  Two normalized pressures for the same velocity differ by an
\(L^2\) function with zero gradient and hence coincide.  The uniqueness
already proved is consequently uniqueness of the pair \((w,q)\).

\divider{Higher regularity without a shorter clock}

Fix an integer \(m\geq3\) and temporarily take \(a\) smooth.  Differentiate the
projected equation by every \(\partial^\alpha\), \(|\alpha|\leq m\), pair with
the same derivative of \(w\), and sum.  Divergence-free transport cancels at the
top order:
\[
  \bigl((w\cdot\nabla)\partial^\alpha w,\partial^\alpha w\bigr)_2=0.
\]
Every remaining Leibniz term differentiates the transporting factor at least
once.  The embeddings \(H^2\hookrightarrow L^\infty\) and
\(H^1\hookrightarrow L^6\) then give
\[
  \left(
  \sum_{|\alpha|\leq m}
  \norm{[\partial^\alpha,w\cdot\nabla]w}_2^2
  \right)^{1/2}
  \leq C_m\norm{\nabla w}_\infty\norm{w}_{H^m}.
\]
Thus the differentiated energy obeys
\begin{equation}
  \frac12\frac d{dt}\norm{w}_{H^m}^2
  +\nu\norm{\nabla w}_{H^m}^2
  \leq C_m\norm{\nabla w}_\infty\norm{w}_{H^m}^2.
  \label{eq:LocalPersistenceEnergy}
\end{equation}
The already constructed \(C_tH^3_x\) solution has integrable Lipschitz norm on
this same interval.  Gronwall yields
\begin{equation}
  \norm{w(t)}_{H^m}^2
  \leq
  \norm{a}_{H^m}^2
  \exp\left(
    2C_m\int_{t_0}^{t}\norm{\nabla w(s)}_\infty\,ds
  \right).
  \label{eq:LocalPersistenceBound}
\end{equation}
Retaining the dissipative term gives \(w\in L^2(J_\delta;H^{m+1})\).  At order
\(m-1\), the nonlinear term lies in \(L^2(J_\delta;H^{m-1})\), so the equation
also gives
\[
  w_t\in L^2(J_\delta;H^{m-1}).
\]
The adjacent trace at height \(m\) completes the three rows:
\begin{equation}
  w\in C(J_\delta;H^m)
  \cap L^2(J_\delta;H^{m+1}),
  \qquad
  w_t\in L^2(J_\delta;H^{m-1}).
  \label{eq:LocalHigherPersistence}
\end{equation}
The pressure formula and the \(H^m\) algebra property put
\(q\) in \(C(J_\delta;H^m)\) as well.

For general \(a\in H^m_\sigma\), approximate it by divergence-free frequency
cutoffs \(a_k\).  Their \(H^3\) norms share a bound controlled by
\(\norm{a}_{H^3}\); hence \eqref{eq:LocalLifespan} puts every smooth
approximation on one common interval.  The fixed-point estimate identifies the
\(H^3\) limit with \(w\), while the persistence inequalities bound all three
higher rows uniformly.  Weak compactness then identifies their limits with the
same \(w\).  The adjacent trace gives continuity in \(H^m\); the energy
inequality at \(t_0\), together with weak convergence to \(a\), upgrades the
initial convergence to strong convergence in \(H^m\).

The interval never depends on \(m\); it was chosen once from \(\nu\) and
\(\norm{a}_{H^3}\).  If \(a\in H^\infty_\sigma\), repeat the argument at each
finite height.  Uniqueness in \(H^3\) keeps every height on the same pair, and
the pressure formula and momentum equation then recover every finite time
derivative.  Spatial embedding makes \((w,q)\) smooth throughout the local
space--time slab.

\divider{Joining the local pieces and locating the only obstruction}

Begin at \(t_0\) and extend through every interval allowed by the construction.
Strong uniqueness makes the velocity and normalized pressure identical on
overlaps, producing one pressure-complete history on a maximal interval
\([t_0,T_{\max})\).

Suppose \(T_{\max}<\infty\) and
\[
  \sup_{t_0\leq t<T_{\max}}\norm{w(t)}_{H^3}\leq M.
\]
If its displayed \(H^3\) bound were finite, the lifespan estimate would supply
the same positive amount of future time
\[
  \delta_M
  :=\min\left\{1,c_\nu(1+M)^{-2}\right\}
\]
from every later datum \(w(t)\).  Choose
\(t>T_{\max}-\delta_M/2\).  Restarting at \(w(t)\) then crosses
\(T_{\max}\), and uniqueness glues that restart to the supposed maximal
history.  The contradiction proves
\begin{equation}
  T_{\max}<\infty
  \quad\Longrightarrow\quad
  \limsup_{t\uparrow T_{\max}}\norm{w(t)}_{H^3}=\infty.
  \label{eq:LocalContinuationAlternative}
\end{equation}
Every higher order already present in the datum persists along each piece,
because every piece uses this same \(H^3\)-controlled clock.  In particular,
\(a\in\Ssigma\subset H^\infty_\sigma\) generates the smooth normalized maximal
history claimed above.
\end{proof}
